% TOP_PROBLEM: {"problemId": "problem.top500-omission-w041044-weight-monodromy-conjecture", "canonicalTitle": "Weight–Monodromy Conjecture", "releaseRank": 85, "primaryDomain": "number_theory_arithmetic_geometry", "primaryDomainLabel": "Number theory & arithmetic geometry", "displayStatus": "open_with_solved_subcases", "releaseStatus": "open_with_solved_subcases"}
% !TeX program = lualatex
% Definition and sources only; this file is not an LLM solution attempt.
\documentclass[11pt]{article}
\usepackage[margin=25mm]{geometry}
\usepackage{fontspec}
\setmainfont{Latin Modern Roman}
\usepackage{amsmath,amssymb,mathtools}
\usepackage{unicode-math}
\setmathfont{Latin Modern Math}
\usepackage{xurl,hyperref}
\hypersetup{colorlinks=true,urlcolor=blue}
\setcounter{secnumdepth}{0}
\setlength{\parindent}{0pt}
\setlength{\parskip}{5pt}
\setlength{\emergencystretch}{3em}
\begin{document}
\section*{\texorpdfstring{Weight–Monodromy Conjecture}{Weight–Monodromy Conjecture}}\label{85}
\subsection{Definitions and mathematical statement}
Put $V=H^i_{\mathrm{\acute et}}(X_{\bar K},\overline{\mathbb{Q}}_\ell)$ for a proper smooth variety over a nonarchimedean local field with residue cardinality $q$, and $\ell\ne p$. Inertia defines nilpotent monodromy $N:V\to V(-1)$. Its zero-centered increasing filtration $M$ is characterized by
\[
 N(M_j)\subseteq M_{j-2}(-1),\qquad
 N^a:\operatorname{gr}^M_aV\xrightarrow{\sim}\operatorname{gr}^M_{-a}V(-a).
\]
For any geometric Frobenius lift, every eigenvalue $\alpha$ on $\operatorname{gr}^M_jV$ should be algebraic and satisfy $|\iota(\alpha)|=q^{(i+j)/2}$ for every complex embedding $\iota$. This is purity of weight $i+j$. The original field and its $q$ remain fixed when defining monodromy on an open inertia subgroup; no semistability or projectivity assumption is added.
\par\smallskip\textit{Formulation and concept references:} \href{https://arxiv.org/pdf/1111.4914v1}{[S1]}.

\subsection{Short English statement}
Do monodromy layers of local-field étale cohomology have exactly the Frobenius weights predicted by their degrees?
\subsection{Sources}
\begin{itemize}
\item[S1]Peter Scholze, Perfectoid spaces, arXiv1111.4914v1,banner21November2011/internal27November2024,51pages.
\url{https://arxiv.org/pdf/1111.4914v1}.
\textit{Formulation source.}
\item[S2]Formality for rigid-analytic spaces satisfying the weight-monodromy conjecture.
\url{https://arxiv.org/abs/2607.14517}.
\textit{Further reference; not a proof endorsement.}
\item[S3]Perfectoid covers of abelian varieties and the weight-monodromy conjecture.
\url{https://arxiv.org/abs/2303.05610}.
\textit{Further reference; not a proof endorsement.}
\item[S4]Semistable Lefschetz pencils.
\url{https://page.mi.fu-berlin.de/esnault/preprints/helene/151_mw.pdf}.
\textit{Further reference; not a proof endorsement.}
\end{itemize}
\end{document}
